\documentclass[11pt]{scrartcl}
\usepackage[sexy]{evan}
\usepackage{listings}
\usepackage{booktabs}

\begin{document}
\title{Enough Python to solve $N = p^2 + q^3$}
\author{Andrew Lin}
\date{September 2026}
\maketitle

My goal is to explain enough Python to solve the programming problem in the OTIS application.
Some statements here are technically false for the purpose of simplifying the exposition.
Sample programs are given for each section.

\section{Variables}
A variable is a name assigned to a value.
To create or change a variable, put the variable name before an = sign,
then put the value (which can include variables already created). Variables can be the following:
\begin{itemize}
    \item A real number, written in decimal form.
    \item A comma-separated list of numbers bounded by [ ].
\end{itemize}
The value of a variable will not change unless you change it directly.
\begin{lstlisting}[language=Python]
x = 3
y = [-1/4, 3/4]
z = x # z is 3
x = 4 # z is still 3
\end{lstlisting}
Some words are \emph{keywords} with a special purpose and cannot be used as the name of a variable.
When using a keyword for its special purpose,
there must be a space between the keyword and any letter or number.

\section{Math operations}
Python supports the following operations:
\begin{center}
\begin{tabular}{cc}
    \toprule
    Symbol & Operation, when placed between two numbers \\ \midrule
    + & addition \\
    - & subtraction \\
    * & multiplication \\
    / & division \\
    ** & exponentiation \\
    \% & modulo \\
    \bottomrule
\end{tabular}
\end{center}
The + sign, when placed between two lists of numbers, concatenates them.
\begin{lstlisting}[language=Python]
x = 10
y = 3
a = x + y # a is 13
b = x - y # b is 7
c = x * y # c is 30
d = x / y # d is 10/3
e = x ** y # e is 1000
f = x % y # f is 10 (mod 3) = 1
z = [1, 2, 3] + [4, 5] # z is [1, 2, 3, 4, 5]
\end{lstlisting}

\section{Functions}
Like in math, functions produce some output given some inputs.
Define a function using the \textbf{def} keyword, followed by the function name,
a comma-separated list of parameters bounded by parentheses, and a colon.

The following lines should be an indented block of code that determines the output given the parameters.
To specify the output, use the \textbf{return} keyword followed by the data that should be outputted.
The function can now be used by using the function name followed by a separated
list of values bounded by parentheses (your inputs).

Once the program reads a single return statement, the rest of the function is not read.
Any variables defined in a function cannot be accessed by the rest of the program.
\begin{lstlisting}[language=Python]
def double(a):
    b = 2 * a
    c = 1434
    return b
    return 665
z = double(5) # z is 10
y = c # invalid
\end{lstlisting}
There are some \emph{built-in} functions.
The built-in \textbf{print} function takes one input and outputs it to the screen so you can see it.
\begin{lstlisting}[language=Python]
x = 264
print(x ** 2) # your screen will show 69696
\end{lstlisting}

\section{Conditionals}
An if statement allows you to run a block of code only if a condition is satisfied.
Use the \textbf{if} keyword, followed by a condition and a colon.
The following lines should be an indented block of code that runs if and only if the condition is satisfied.
Optionally, immediately after this indented block, use the \textbf{else} keyword, followed by a colon.
Again, the following lines should be an indented block of code.
This block of code will run only if the original condition does not hold.
\subsection{Conditions}
The following symbols can be used to create conditions:
\begin{center}
\begin{tabular}{cc}
    \toprule
    Symbol & Meaning \\ \midrule
    == & equal to \\
    $>$ & greater than \\
    $<$ & less than \\
    $>=$ & greater than or equal to \\
    $<=$ & less than or equal to \\
    != & not equal to \\
    \bottomrule
\end{tabular}
\end{center}
Python supports the following logical operators:
\begin{itemize}
    \item The \textbf{not} keyword, placed in front of a condition $C$, creates a condition which holds if and only if $C$ does not hold.
    \item The \textbf{and} keyword, placed between two conditions $C$ and $C'$, creates a condition that holds if and only if both $C$ and $C'$ hold.
    \item The \textbf{or} keyword, placed between two conditions $C$ and $C'$, creates a condition that holds if and only if at least one of $C$ and $C'$ holds.
\end{itemize}

\begin{lstlisting}[language=Python]
x = 7
if x > 0 and x % 2 == 0:
    x = x / 2
else:
    x = 3 * x + 1
# only the second body runs, so x is now 22
\end{lstlisting}

\section{Loops}
Loops allow you to run the same piece of code multiple times.
\subsection{For loops}
Use the \textbf{for} keyword, followed by a name used nowhere else in the program, the \textbf{in} keyword,
a list of numbers, and a colon. The following lines should contain an indented block of code.
The program iterates over the list of numbers,
setting a variable with the specified name equal to the number in the list, and running the code block.

The built-in \textbf{range} function takes as input two integer inputs $a < b$
and outputs a list of all integers in order from $a$ to $b$, including $a$ but not $b$.
\begin{lstlisting}[language=Python]
a = 0
for i in range(1, 11):
    a = a + i
# range(1, 11) is [1, 2, ..., 10]; for each number in turn we set i equal to the number and add it to a
# a = 1 + 2 + ... + 10 = 55
\end{lstlisting}

\subsection{While loops}
Use the \textbf{while} keyword, followed by a condition and a colon.
The following lines should contain an indented block of code.
The program runs the code in the block until the condition does not hold;
it is possible that the code is never run at all.
\begin{lstlisting}[language=Python]
a = 7
while a != 1:
    if a % 2 == 0:
        a = a / 2
    else:
        a = 3 * a + 1
# runs the Collatz process starting with 7 until we reach 1
\end{lstlisting}
\end{document}
